\documentclass[11pt,a4paper]{article}
\usepackage[utf8]{inputenc}
\usepackage[T1]{fontenc}
\usepackage[english]{babel}
\usepackage{amsmath, amssymb, amsthm, mathtools}
\usepackage{geometry}
\usepackage{hyperref}
\usepackage{csquotes}
\usepackage{enumitem}
\usepackage{tikz-cd}
\usepackage[most]{tcolorbox}

\geometry{margin=2.5cm}
\hypersetup{
	colorlinks=true,
	linkcolor=black,
	citecolor=black,
	urlcolor=blue
}

\title{Mathematics Through Symmetry: \\ Orbit Groupoids, Twin Structures, and a Universal \\ Construction and Classification Theorem}
\author{Jocelyn Nembe \\ Modeling and Calculus Laboratory, ROOTS INSIGHTS \\ Libreville, Gabon}
\date{December 2025}

\newtheorem{theorem}{Theorem}[section]
\newtheorem{definition}[theorem]{Definition}
\newtheorem{proposition}[theorem]{Proposition}
\newtheorem{corollary}[theorem]{Corollary}
\newtheorem{lemma}[theorem]{Lemma}
\newtheorem{remark}[theorem]{Remark}
\newtheorem{example}[theorem]{Example}

\begin{document}
	
	\maketitle
	
	\begin{abstract}
		We develop a universal mathematical framework for constructing and classifying structural deformations induced by group symmetries. Given a binary structure $(E, \diamond)$ and a group $G$ acting on $E$, our procedure canonically generates new operations $\diamond_g$ parameterized by group elements. We prove a universal classification theorem showing that the deformation space is canonically isomorphic to the homogeneous space $L\backslash G$ (right cosets), where $L$ is the structural kernel preserving the original operation. While the individual constructions (e.g.\ pullback of metrics, conjugation of operators) are classical, our contribution lies in unifying them under a single orbit-theoretic principle. Through purely algebraic and combinatorial examples, we illustrate how this perspective reveals a common mechanism behind apparent diversity in algebra and geometry. The framework does not introduce new physics or solve open problems in analysis; rather, it offers a conceptual lens for classifying symmetry-induced structural change.
	\end{abstract}
	
	\section{Introduction}
	
	Symmetries play a central role in the foundational organization of mathematics \cite{Klein1872,Ehresmann1950}. In category theory, topology, geometry, and algebraic structures, group actions express invariance and govern deformations \cite{Bourbaki,MacLane}. In this work, we demonstrate that the essence of structural transformation under symmetry can be captured by a general \emph{transport operation}: a single mechanism that both constructs deformations and classifies them.
	
	Our main contributions are:
	\begin{itemize}
		\item A universal construction generating all group-induced deformations of $(E,\diamond)$;
		\item A classification theorem proving that all deformations are parametrized by $L\backslash G$, where $L$ is the structural kernel;
		\item Purely mathematical examples—from deformed group laws to twisted convolutions—illustrating the breadth of the formalism.
	\end{itemize}
	
	This establishes a unified principle applicable across algebra, combinatorics, and geometry.
	
	\subsection*{What this work is \emph{not}}
	This article does \emph{not}:
	\begin{itemize}
		\item propose a new physical theory unifying gravity and quantum mechanics;
		\item introduce novel analytical theorems in differential geometry or operator theory;
		\item claim technical novelty for constructions like metric pullback or operator conjugation.
	\end{itemize}
	Rather, it offers a \emph{structural reinterpretation}: it shows that these classical constructions are manifestations of a single algebraic transport mechanism, and that their space of deformations is universally classified by a homogeneous space $L\backslash G$.
	
	\subsection*{Symmetry generates both rigidity and flexibility}
	
	Some structures are rigid under symmetries (e.g., classical commutativity, flat geometry), while others exhibit flexibility (curvature, noncommutativity). We show that both phenomena result from whether $G = L$ or $[G : L] > 1$.
	
	\subsection*{Construct \emph{and} classify}
	
	A key aspect of our contribution is that the framework:
	\begin{itemize}
		\item constructs all possible symmetry-driven deformations; and
		\item classifies them completely using the single invariant $L\backslash G$.
	\end{itemize}
	No additional data is needed: the action of $G$ on $(E, \diamond)$ determines the entire space of outcomes.
	
	\section{Orbit Classification by Group Actions}
	
	A group action $G \curvearrowright E$ induces an equivalence relation: $x \sim y \iff \exists g \in G : y = g \cdot x$.
	The equivalence classes $G \cdot x$ are the \emph{orbits} \cite{Serre}.
	
	\begin{theorem}[Orbit Partition Theorem]
		$E$ decomposes into disjoint orbits:
		\[
		E = \bigsqcup_{x \in E/G} G \cdot x.
		\]
	\end{theorem}
	
	\begin{definition}[Stabilizer]
		$G_x := \{ g \in G : g \cdot x = x \}$.
	\end{definition}
	
	\begin{theorem}[Orbit–Stabilizer Theorem]
		There is a $G$-equivariant bijection
		\[
		G/G_x \simeq G \cdot x.
		\]
	\end{theorem}
	
	\begin{remark}
		Homogeneous spaces $G/G_x$ classify configurations of $x$ generated by symmetries.
	\end{remark}
	
	\section{Preliminaries}
	
	Let $G$ be a group acting on a nonempty set $E$ by $G \times E \to E$, $(g,x) \mapsto g \cdot x$.
	We fix throughout a binary operation $\diamond : E \times E \to E$.
	
	Transport of the structure by $g \in G$ is defined by
	\begin{equation}
		x \diamond_g y := g^{-1} \cdot \big((g \cdot x) \diamond (g \cdot y)\big), \label{eq:transport}
	\end{equation}
	which yields a new binary structure $(E, \diamond_g)$.
	
	\begin{remark}
		This construction generalizes conjugation in groups and twisted products in noncommutative geometry \cite{Connes}.
	\end{remark}
	
	Denote the set of transported operations by $\mathcal{O} := \{ \diamond_g : g \in G \}$.
	
	\begin{tcolorbox}[title=Box 1: Operations vs.~Objects – An Inversion of Priority, colback=gray!5!white, colframe=gray!75!black]
		Traditionally, one transports \emph{objects} (functions, tensors, operators) under group actions, keeping the algebraic structure fixed. Here, we fix the underlying set $E$ and \emph{transport the operation itself}. This inversion reveals that what varies is not the carrier set, but the \emph{rule of composition}. The formalism $x \diamond_g y := g^{-1}((g\cdot x) \diamond (g\cdot y))$ makes this inversion explicit and universal.
		
		\medskip
		\textbf{Concrete illustration}: On $\mathbb{R}$ with addition, if $g: x \mapsto x + c$, the transported operation becomes $x \oplus_g y = x + y + c$. The carrier set $\mathbb{R}$ is unchanged; only the rule of combination has shifted.
	\end{tcolorbox}
	
	\section{Twin Spaces}
	
	We call
	\[
	g \cdot (E, \diamond) := (E, \diamond_g)
	\]
	a \emph{twin space}. Thus $G$ acts on the set of all twin spaces, forming a groupoid.

More precisely, the relevant object is the \emph{transport (action) groupoid} $G\ltimes\mathcal O$: its objects
are the twin spaces $\diamond_g\in\mathcal O$, and a morphism $\diamond_g\to\diamond_{g'}$ is an element $h\in G$
carrying the first to the second. Since every twin space is reached from $\diamond$, this groupoid is connected,
and its isotropy (vertex) group at the base point $\diamond$ is exactly the structural kernel $L$. A connected
groupoid is equivalent to its vertex group, so $G\ltimes\mathcal O\simeq L$ as a groupoid, while its set of
objects is the homogeneous space $\mathcal O\cong L\backslash G$ of the classification theorem below. The
groupoid therefore packages both halves of the theory at once: its objects record the \emph{deformations}
$L\backslash G$, and its vertex group records the \emph{symmetries} $L$.
	
	\begin{lemma}
		The assignment $g \mapsto g \cdot (E, \diamond)$ defines a transitive right action of $G$ on the set of twin spaces (transport composes on the right: $(\diamond_g)_h = \diamond_{gh}$).
	\end{lemma}
	\begin{proof}
		Immediate from \eqref{eq:transport}.
	\end{proof}
	
	Twin structures inherit many algebraic and geometric properties of the original structure, but often with transformed parameters.
	
	\section{Twin Operations and the Structural Kernel}
	
	We now introduce the map that associates to each $g \in G$ its transported operation $\diamond_g$.
	This defines the \emph{deformation morphism} $\pi : G \to \mathcal{O}$, $\pi(g) = \diamond_g$.
	
	\begin{definition}[Structural kernel]
		The structural kernel is
		\[
		L := \{ g \in G : \diamond_g = \diamond \} = \{ g \in G : g \cdot (x \diamond y) = (g \cdot x) \diamond (g \cdot y) \}.
		\]
	\end{definition}
	
	\begin{lemma}
		$L$ is a subgroup of $G$.
	\end{lemma}
	\begin{proof}
		The map $g\mapsto\diamond_g$ is an action of $G$ on $\mathcal O$ and $L$ is the stabiliser of the base point
		$\diamond$, so it is a subgroup. Explicitly $\diamond_e=\diamond$ gives $e\in L$; if $\diamond_g=\diamond$ then
		applying the transport by $g^{-1}$ gives $\diamond_{g^{-1}}=\diamond$; and if $\diamond_g=\diamond_h=\diamond$
		then $\diamond_{gh}=\diamond$ by the composition law of \eqref{eq:transport}.
	\end{proof}
	
	\begin{remark}
		Elements of $L$ preserve the original structure. Thus $L$ encodes the internal symmetry of $(E, \diamond)$.
	\end{remark}
	
	\section{Universal Construction and Classification Theorem}
	
	We now state our main result, showing that the transported structures form a homogeneous space.
	
	\begin{theorem}[Universal Construction and Classification Theorem]
		There is a canonical bijection $\mathcal{O} \simeq L\backslash G$, given by $Lg \mapsto \diamond_g$.
		The fibers of $\pi$ are the right cosets $Lg$; thus $\mathcal{O} \cong L\backslash G$ as a homogeneous space for the right transport action.
	\end{theorem}
	
	\begin{proof}
		Surjectivity of $\pi$ follows by definition of $\mathcal{O}$.
		
		Let $g_1, g_2 \in G$. Then $\diamond_{g_1} = \diamond_{g_2}$ iff
		\[
		g_2 g_1^{-1} \cdot (x \diamond y) = (g_2 g_1^{-1} \cdot x) \diamond (g_2 g_1^{-1} \cdot y) \quad \forall x,y \in E,
		\]
		i.e.\ iff $g_2 g_1^{-1} \in L$. Thus $\diamond_{g_1} = \diamond_{g_2}$ exactly when $g_1,g_2$ lie in the same right coset (indeed $g_2 g_1^{-1} \in L \iff L g_1 = L g_2$).
		
		The fibers of $\pi$ are thus precisely the right cosets $Lg$. Hence $\mathcal{O} \simeq L\backslash G$.
	\end{proof}
	
	\begin{corollary}[Universality]
		Every transported structure in $\mathcal{O}$ arises uniquely from some $g \in G$, modulo the structural kernel $L$:
		\[
		\mathcal{O} = \{ \diamond_g : g \in G \} \quad \text{and} \quad \mathcal{O} \simeq L\backslash G.
		\]
		No additional data is needed to generate or parametrize deformations.
	\end{corollary}
	
	\begin{corollary}[Rigidity versus flexibility]
		\[
		G = L \iff \mathcal{O} = \{ \diamond \}.
		\]
		Nontrivial transported structures exist iff $[G : L] > 1$.
	\end{corollary}
	
	\begin{remark}
		The deformation space depends only on the action of $G$ and not on the specific form of $\diamond$ beyond its symmetry group.
	\end{remark}

	\begin{remark}[Right cosets, and normality]
		The fibers of $\pi$ are the \emph{right} cosets $Lg$ (the condition is $g_2 g_1^{-1} \in L$); the canonical model is therefore $L\backslash G$, not $G/L$. Inversion $g \mapsto g^{-1}$ is a bijection $L\backslash G \to G/L$, so the \emph{number} of deformations is the index $[G:L]$ either way, but the natural assignment $g \mapsto \diamond_g$ is constant precisely on right cosets. The two coset spaces coincide canonically iff $L \trianglelefteq G$. In the affine examples below, $L$ is not normal, and the deformation parameter is exactly a right-coset invariant.
	\end{remark}
	
	\section{Purely Mathematical Examples}
	
	We now illustrate the universality of the transport mechanism through \emph{purely mathematical} (non-physical) examples.
	
	\begin{example}[Translated addition on $\mathbb{Z}$]
		Let $E = \mathbb{Z}$ with $\diamond = +$. Let $G = \mathrm{Aff}(\mathbb{Z}) = \{ x \mapsto \varepsilon x + b \mid \varepsilon = \pm 1, b \in \mathbb{Z} \}$.
		For $g(x) = x + 1$, we obtain
		\[
		x \oplus_g y = g^{-1}(g(x) + g(y)) = ((x+1) + (y+1)) - 1 = x + y + 1.
		\]
		This operation remains \emph{associative}: $(x \oplus_g y) \oplus_g z = x + y + z + 2 = x \oplus_g (y \oplus_g z)$.
		In fact, $(E, \oplus_g)$ is isomorphic to $(E, +)$ via the translation $x \mapsto x + 1$.
		
		Here $L = \{ x \mapsto \varepsilon x \mid \varepsilon = \pm 1 \}$, and $L\backslash G \cong \mathbb{Z}$ parametrizes a family of ``translated'' additions, all mutually isomorphic but formally distinct as operations on $\mathbb{Z}$.
	\end{example}
	
	\begin{example}[Twisted convolution algebras]
		Let $E = L^1(G)$ with convolution $\ast$. For $\alpha \in \mathrm{Aut}(G)$, define
		\[
		f \ast_\alpha h := \alpha^{-1}(\alpha(f) \ast \alpha(h)).
		\]
		Because convolution is equivariant under \emph{every} automorphism, $\alpha(f) \ast \alpha(h) = \alpha(f \ast h)$ (up to the modular factor on non-unimodular $G$), so $f \ast_\alpha h = f \ast h$ for all $\alpha$. Hence the structural kernel is $L = \mathrm{Aut}(G)$ and the deformation is \emph{trivial}: convolution is rigid under automorphism transport. Genuinely ``twisted'' convolution algebras instead arise from a $2$-cocycle $\sigma \in Z^2(G,\mathbb{T})$ (the twisted group algebra $\mathbb{C}_\sigma[G]$) --- a \emph{cohomological} deformation, not a transport, and so outside the present mechanism.
	\end{example}
	
	\begin{example}[Deformed lattice operations]
		Let $E = \mathcal{P}(X)$ with intersection $\cap$. Let $G = \mathrm{Sym}(X)$ act by permutation: $\sigma \cdot A = \sigma(A)$.
		Then $A \cap_\sigma B = \sigma^{-1}(\sigma(A) \cap \sigma(B)) = A \cap B$, so $L = G$: the lattice is \emph{rigid}. This illustrates that set-theoretic operations often exhibit maximal rigidity under natural symmetry groups.
	\end{example}
	
	\section{Four Explicit Computations of $L\backslash G$}
	
	We now present detailed computations of the deformation space for various structures.
	
	\subsection{Example 1: Deformed addition on $\mathbb{R}$ under affine transformations}
	
	\textbf{Setup:}
	\begin{itemize}[nosep]
		\item Set: $E = \mathbb{R}$
		\item Operation: $x \diamond y = x + y$
		\item Group: $G = \mathrm{Aff}^+(\mathbb{R}) = \{ g_{a,b}: x \mapsto ax + b \mid a > 0,\, b \in \mathbb{R} \}$
	\end{itemize}
	
	\textbf{Computation:} For $g = g_{a,b}$, we have $g^{-1}(x) = (x-b)/a$, so
	\[
	x \diamond_g y = g^{-1}(g(x) + g(y)) = \frac{(ax+b) + (ay+b) - b}{a} = x + y + \frac{b}{a}.
	\]
	
	\textbf{Structural kernel:} $\diamond_g = \diamond$ iff $b/a = 0$ iff $b = 0$. Thus
	\[
	L = \{ x \mapsto ax \mid a > 0 \} \cong (\mathbb{R}_{>0}, \times).
	\]
	
	\textbf{Classification:} $L\backslash G \cong \mathbb{R}$, with the right coset $L g_{a,b}$ determined by $b/a \in \mathbb{R}$. Note $L$ is \emph{not} normal in $G$, so $L\backslash G \neq G/L$ canonically; $b/a$ is the right-coset invariant ($s_\alpha g_{a,b} = g_{\alpha a, \alpha b}$ fixes it), whereas the left coset $g_{a,b}L$ is labelled by $b$.
	
	The transported operations $x \oplus_c y = x + y + c$ for $c \in \mathbb{R}$ form a one-parameter family of group laws on $\mathbb{R}$, all isomorphic to $(\mathbb{R}, +)$.
	
	\subsection{Example 2: Conjugation action on matrix multiplication (Rigidity)}
	
	\textbf{Setup:}
	\begin{itemize}[nosep]
		\item Set: $E = \mathrm{M}_n(\mathbb{C})$
		\item Operation: $A \diamond B = AB$ (matrix multiplication)
		\item Group: $G = \mathrm{GL}_n(\mathbb{C})$, acting by conjugation: $g \cdot A = gAg^{-1}$
	\end{itemize}
	
	\textbf{Computation:}
	\begin{align*}
		A \diamond_g B &= g^{-1} \cdot \big((g \cdot A) \diamond (g \cdot B)\big) \\
		&= g^{-1} \cdot \big((gAg^{-1})(gBg^{-1})\big) \\
		&= g^{-1} \cdot (gABg^{-1}) \\
		&= g^{-1}(gABg^{-1})g = AB = A \diamond B.
	\end{align*}
	
	\textbf{Structural kernel:} $L = G$ (every element preserves the operation).
	
	\textbf{Classification:} $L\backslash G = \{*\}$ (single point). The deformation space is trivial: $\mathcal{O} = \{\diamond\}$.
	
	\textbf{Interpretation:} Matrix multiplication exhibits \emph{complete structural rigidity} under conjugation. This reflects the fact that conjugation is an inner automorphism of the algebra $\mathrm{M}_n(\mathbb{C})$.
	
	\subsection{Example 3: Cyclic group acting on a finite magma}
	
	\textbf{Setup:}
	\begin{itemize}[nosep]
		\item Set: $E = \{0, 1, 2\}$
		\item Operation: Define the magma $(E, \diamond)$ by the following Cayley table:
	\end{itemize}
	
	\begin{center}
		\begin{tabular}{c|ccc}
			$\diamond$ & 0 & 1 & 2 \\ \hline
			0 & 0 & 1 & 2 \\
			1 & 2 & 0 & 1 \\
			2 & 1 & 2 & 0
		\end{tabular}
	\end{center}
	
	\begin{itemize}[nosep]
		\item Group: $G = C_3 = \langle \sigma \rangle$ where $\sigma = (0\;1\;2)$ acts by cyclic permutation.
	\end{itemize}
	
	\textbf{Computation:} The transported table under $\sigma$ is:
	\[
	x \diamond_\sigma y = \sigma^{-1}(\sigma(x) \diamond \sigma(y)).
	\]
	For instance, $0 \diamond_\sigma 1 = \sigma^{-1}(1 \diamond 2) = \sigma^{-1}(1) = 0$.
	
	Computing the full table:
	\begin{center}
		\begin{tabular}{c|ccc}
			$\diamond_\sigma$ & 0 & 1 & 2 \\ \hline
			0 & 2 & 0 & 1 \\
			1 & 1 & 2 & 0 \\
			2 & 0 & 1 & 2
		\end{tabular}
	\end{center}
	
	This is distinct from $(E, \diamond)$ since $1 \diamond 0 = 2 \neq 1 = 1 \diamond_\sigma 0$.
	
	\textbf{Structural kernel:} Direct verification shows $L = \{e\}$ (only the identity preserves the table).
	
	\textbf{Classification:} $L\backslash G \cong C_3$, yielding three distinct operations $\diamond$, $\diamond_\sigma$, $\diamond_{\sigma^2}$.
	
	\subsection{Example 4: Two actions on $(\mathbb{R}_{>0}, \cdot)$}
	
	This example shows how the \emph{same operation} can exhibit different rigidity depending on the group action.
	
	\textbf{Setup:} $E = \mathbb{R}_{>0}$, $\diamond = \cdot$ (multiplication).
	
	\textbf{Action 1: Multiplicative scaling.} Let $G = (\mathbb{R}, +)$ act by $t \cdot x = e^t x$.
	
	Then $x \diamond_t y = e^{-t}(e^t x \cdot e^t y) = e^t xy \neq xy$ for $t \neq 0$.
	
	Structural kernel: $L = \{0\}$. Quotient: $L\backslash G = \mathbb{R}$ (infinite deformation space).
	
	\textbf{Action 2: Power action.} Let $G = (\mathbb{R}_{>0}, \times)$ act by $a \cdot x = x^a$.
	
	Then $x \diamond_a y = (x^a \cdot y^a)^{1/a} = (xy)^{a \cdot (1/a)} = xy$.
	
	Structural kernel: $L = G$. Quotient: $L\backslash G = \{*\}$ (complete rigidity).
	
	\textbf{Conclusion:} The deformation space depends crucially on \emph{how} the group acts, not just on the group or the operation individually.
	
	\section{Conceptual Contribution}
	
	While the orbit-stabilizer theorem is classical, its systematic application to \emph{operation deformation}—rather than point orbits—appears underexploited in the literature. Our framework makes this application explicit and uniform.
	
	\subsection*{A unified viewpoint}
	
	The transport formula $\diamond_g = g^{-1} \circ \diamond \circ (g \times g)$ reveals that:
	\begin{itemize}
		\item Metric pullback under diffeomorphisms (geometry),
		\item Product deformation under automorphisms (algebra),
		\item Conjugation of group laws (group theory)
	\end{itemize}
	are instances of a \emph{single mechanism}. This is not merely a metaphor but a precise structural identity.
	
	\subsection*{Classification without cohomology}
	
	In some contexts, deformation classification requires cohomological machinery (e.g., $H^2(G, A)$ for group extensions). The present framework provides an alternative when deformations arise purely from transport: the classification is immediate from $L\backslash G$ without spectral sequences or cocycle computations. The two
viewpoints are complementary rather than opposed: the global orbit $L\backslash G$ is the set of finite
deformations of transport type, and its infinitesimal structure at the base point---the tangent space
$\mathfrak g/\mathfrak l$ of the homogeneous space---sits inside the first-order deformation space
$H^1(G,\mathrm{End}\,E)$ classified cohomologically in the companion deformation paper. Transport deformations
are the integrable, transport-type directions; genuinely cohomological deformations (cocycles that are not
transports, such as the $2$-cocycle twists above) lie in the complement. Conversely, genuinely cohomological deformations --- such as $2$-cocycle (projective) twists, as in the convolution example above --- are \emph{not} of transport type and lie outside this framework.
	
	\subsection*{Distinguishing rigidity from flexibility}
	
	The dichotomy $G = L$ (rigidity) versus $[G:L] > 1$ (flexibility) provides a simple criterion for determining whether a structure admits nontrivial deformations under a given symmetry group.
	
	\section{Categorical Perspective}
	
	\begin{proposition}[Universal property]
		Let $\mathcal{C}$ be a category with a forgetful functor to $\mathbf{Set}$. Let $(E,\diamond) \in \mathcal{C}$, and let $G$ act on the underlying set. Then the transport $\diamond_g$ is the unique operation making the diagram
		\[
		\begin{tikzcd}
			E \times E \arrow[r, "\diamond"] \arrow[d, "g^{-1} \times g^{-1}"'] & E \arrow[d, "g^{-1}"] \\
			E \times E \arrow[r, "\diamond_g"'] & E
		\end{tikzcd}
		\]
		commute in $\mathbf{Set}$. Thus $\diamond_g = g^{-1} \circ \diamond \circ (g \times g)$ is the canonical equivariant lift.
	\end{proposition}
	
	\begin{proof}
		Commutativity of the square reads $g^{-1}\circ\diamond=\diamond_g\circ(g^{-1}\times g^{-1})$ as maps
		$E\times E\to E$. Precomposing with the bijection $g\times g$ gives
		$g^{-1}\circ\diamond\circ(g\times g)=\diamond_g$, which both establishes existence and determines $\diamond_g$
		uniquely, since $g^{-1}$ and $g\times g$ are invertible. This is precisely formula~\eqref{eq:transport}.
	\end{proof}

This shows the transport is not ad hoc, but the unique operation compatible with equivariance.
	
	\section{Conclusion}
	
	We introduced the theory of twin structures: transported operations arising from group symmetries. The following key outcomes were established:
	\begin{enumerate}
		\item A universal method to generate all group-induced deformations of a binary structure $(E, \diamond)$;
		\item A classification theorem proving the canonical identification $\mathcal{O} \simeq L\backslash G$, where $L$ is the structural kernel;
		\item Explicit computations demonstrating the framework across algebra, combinatorics, and analysis.
	\end{enumerate}
	
	This perspective reveals that seemingly disparate phenomena—translated group laws, twisted convolutions, rigid lattices—share a common structural origin. The framework provides a \emph{conceptual unifier} for symmetry-induced structural change.
	
	\subsection*{Future Directions}
	\begin{itemize}
		\item Classification of transported $n$-ary structures;
		\item Topological and metric analysis on the homogeneous space $L\backslash G$;
		\item Connections with formal deformation theory and Hochschild cohomology;
		\item Applications to formal language theory and combinatorial algebra.
	\end{itemize}
	
	\section*{Acknowledgments}
	The author thanks colleagues at the African Institute for Computing Sciences and ROOTS INSIGHTS for discussions on structural mathematics. Special thanks to anonymous reviewers for urging greater conceptual clarity.
	
	\begin{thebibliography}{9}
		\bibitem{Klein1872} F.\ Klein. \emph{Vergleichende Betrachtungen über neuere geometrische Forschungen}. Erlangen, 1872.
		\bibitem{Ehresmann1950} C.\ Ehresmann. Les connexions infinitésimales dans un espace fibré différentiable. \emph{Colloque de Topologie}, Bruxelles, 1950.
		\bibitem{Serre} J.-P.\ Serre. \emph{Linear Representations of Finite Groups}. Springer, 1977.
		\bibitem{Bourbaki} N.\ Bourbaki. \emph{Theory of Sets}. Hermann, 1968.
		\bibitem{MacLane} S.\ Mac Lane. \emph{Categories for the Working Mathematician}. Springer, 1971.
		\bibitem{Connes} A.\ Connes. \emph{Noncommutative Geometry}. Academic Press, 1994.
	\end{thebibliography}
	
\end{document}